\documentclass{beamer}

\usepackage{pb-diagram}
%\usepackage{beamerthemesplit}
\usepackage[french]{babel}
\usepackage[utf8]{inputenc}
%\usetheme{default}
\useoutertheme{infolines}
\usepackage{graphicx, tikz}
\usepackage{comment}



\usepackage[utf8]{inputenc} % clavier et acccent
\usepackage[french]{babel} % langue francaise
\usepackage[french]{nomencl}
\usepackage[T1]{fontenc}
\makenomenclature
\usepackage{geometry} % package pour définir la zone du texte
\usepackage{fancyhdr} % package pour charger le style de la mise en page
%\usepackage{amsfonts} % des symboles
\usepackage{pst-labo}% schémat labo 
\usepackage{amssymb}%  des symboles mathématiques
\usepackage{amsmath} % des symboles mathématiques
\usepackage{mathtools}
\usepackage{graphicx} % pour charger des figures
\usepackage[dvips]{epsfig}% pour charger des figures ".eps"
\usepackage{float}
\usepackage{multicol}  % pour écrire en plusieurs colonnes
\usepackage{pdfpages}
\usepackage{color}  % pour écrire en couleur
%\usepackage{enumitem}
\usepackage{parskip}
\usepackage{amsthm}
\usepackage{hyperref} %pour insérer liens 
\usepackage{ stmaryrd } %pour les ensembles d'entiers
\usepackage{bm} %pour mettre formules en gras
\usepackage{nicematrix}

\usepackage{mathabx}

%\usepackage{mathabx}
\usepackage{dsfont}

\newcommand{\commPL}[1]{{\color{green}{#1}}}
\newcommand{\corrPL}[1]{{\color{blue}{#1}}}
\newcommand{\commCB}[1]{{\color{orange}{#1}}}
\newcommand{\commJB}[1]{{\color{red}{#1}}}

\setcounter{tocdepth}{2}

\graphicspath{{Figures/}}  % toutes les figures doivent être dans le dossier "Figures"
\usepackage{listings}

\usepackage{url}
\usepackage{array}
\usepackage{animate}
\usepackage{appendix}
\usepackage{bm}  % pour mettre des équations en gras
\usepackage[bottom]{footmisc} %JE L'AI RAJOUTE POUR METTRE DES NOTES DE BAS DE PAGE
\usepackage{hyperref} % pour des liens hypertexte

%\usepackage[citestyle=numeric]{biblatex}
%\addbibresource{bibliographie.bib}
%\AtBeginBibliography{\small}

% pour avoir la diagonale inverse dans les matrices
\newcommand{\revdots}{\mathinner
{\mkern1mu\raise1pt\vbox{\kern7pt\hbox{.}}
\mkern2mu\raise4pt\hbox{.}:
\mkern2mu\raise7pt\hbox{.}\mkern1mu}
}


\title[Du Rocq et des pavages]{Apprentissage de Rocq en commençant à formaliser un résultat sur les pavages}

\author[L. BARROCHE]{Loryne BARROCHE}

%\institute[I2M]{Institut de mathématiques de Marseille - Universit\'e d'Aix-Marseille}

\institute[L2 MPCI]{Co-encadrée par Etienne Miquey et Félix Castro}

\date{Juin 2026}



%\newtheorem{thm}{Théorème}

%\newtheorem*{defn}{Définition}
%\newtheorem{prop}{Proposition}
%\newtheorem*{rk}{Remarque}
%\newtheorem*{ex}{Exemple}
%\newtheorem*{exs}{Exemples}
%\newtheorem*{mot}{Motivation}
%\newtheorem*{question}{Question}
%\newtheorem*{cor}{Corollaire}

\newcommand{\N}{\mathbb{N}}
\newcommand{\Z}{\mathbb{Z}}
\newcommand{\R}{\mathbb{R}}
\newcommand{\Q}{\mathbb{Q}}
\newcommand{\restrict}[2]{#1 \upharpoonright #2}
\newcommand{\arrows}[3]{\longrightarrow {#1}^{#2}_{#3}}
\newcommand{\Ur}{\mathbb{U}}
\newcommand{\titre}[1]{\begin{center} \textbf{#1} \end{center}}
\newcommand{\funct}[2]{#1 \longrightarrow #2}
\newcommand{\om}[1]{\textbf{#1},<^{\textbf{#1}}}
\newcommand{\m}[1]{\textbf{#1}}
\newcommand{\oc}[1]{\widetilde{\textbf{#1}},<^{\widetilde{\textbf{#1}}}}
\newcommand{\mc}[1]{\widetilde{\textbf{#1}}}
\newcommand{\iso}{\mathrm{iso}}
\newcommand{\s}{\mathbb{S}}
\newcommand{\dom}{\mathrm{dom}}
\newcommand{\Rt}{\overrightarrow{\mathcal{R}}}
\newcommand{\Cy}{S(2)}
\newcommand{\Aut}{\mathrm{Aut}}

\begin{document}

\begin{frame}
\titlepage

\end{frame}

\begin{frame}
\frametitle{Qu'est-ce que Rocq ?}

\begin{itemize}
    \item Un assistant formalisant des preuves \\
    \item Formaliser une preuve = écrire une preuve de telle sorte à ce que chaque argument découle sans ambiguïté des faits précédents. \\
    Formaliser $\implies$ comprendre tous les détails implicites !
\end{itemize}

\end{frame}


\begin{frame}
\frametitle{Pavage du plan avec des tuiles de Wang}

\begin{itemize}
    \item \textbf{Pavage du plan} = un ensemble de portions du plan recouvrant le plan tout entier sans interstices ni chevauchements\\
    \item \textbf{Tuiles de Wang} = des carrés de longueur 1 définies par la couleur de chacun de ses côtés. Un ensemble fini de tuiles de Wang sera noté $\tau$.\\
    \item \textbf{Pavage du plan avec des tuiles de Wang} = une fonction $P_{2D} : \Z^2 \rightarrow \tau$ qui est "valide" (= compatibilité et recouvrement du plan).
\end{itemize}
\end{frame}


\begin{frame}
\frametitle{Périodicité sur les pavages de Wang}
    \begin{itemize}
        \item  On dit que $P_{2D} : \Z^2 \rightarrow \tau$ est \textbf{faiblement périodique} ssi $\exists \Vec{u} = (x_u, y_u) \neq \Vec{0}$, $(x_u, y_u) \in \Z^2$, $\forall x \in \Z^2$, $\forall k \in \Z$, {$P_{2D}(x+k\Vec{u}) = P_{2D}(x)$}\\
        On dit que $\Vec{u}$ est un vecteur de périodicité de $P_{2D}$.
        \item On dit que $P_{2D} : \Z^2 \rightarrow \tau$ est \textbf{fortement périodique} si et seulement si il existe $\Vec{u} = (x_u, y_u), \Vec{v} = (x_v, y_v)$ deux vecteurs de périodicité de $P_{2D}$ non colinéaires.
    \end{itemize}

\pause

A-t-on fortement périodique $\Leftrightarrow$ faiblement périodique ?\\
Le sens $\Rightarrow$ est évident par définition, mais le sens $\Leftarrow$ l'est beaucoup moins.
\end{frame}




\begin{frame}

    \frametitle{Théorème}
Soit $\tau$ un jeu fini de tuiles de Wang en deux dimensions.\\
$\exists P_{2D} : \Z^2 \rightarrow \tau$ valide, $P_{2D}$ faiblement périodique ssi $\exists P_{2D}^D : \Z^2 \rightarrow \tau$ valide, $P_{2D}^D$ fortement périodique. \\ 

\end{frame}

\begin{frame}
\frametitle{Jeu de tuiles de Wang (non formalisé)}
\NiceMatrixOptions{renew-dots,renew-matrix}
Soit $P_{2D} : \Z^2 \rightarrow \tau$. On forme un \textbf{jeu de tuiles de Wang} de longueur $l+1$ et de hauteur $h+1$ $T_{l+1, h+1}(\tau)$. Il est de la forme : 
$$
T_{l+1, h+1}(\tau) = 
\begin{bmatrix}
t_{h,0} & \Cdots & t_{h,j} & \Cdots & t_{h,l}\\
\Vdots &  & \Vdots &  & \Vdots\\
t_{i, 0} & \Cdots & t_{i,j} & \Cdots & t_{i,l}\\
\Vdots &  & \Vdots &  & \Vdots\\
t_{0,0} & \Cdots & t_{0,j} & \Cdots & t_{0,l}
\end{bmatrix}
$$ \\
On note le pavage (compatibilité + recouvrement) d'une droite par un jeu de tuiles de Wang $T_{l+1, h+1}(\tau)$ avec l'application $P_{1D}^T : \Z \rightarrow T_{l+1, h+1}(\tau)$. 
\end{frame}

\end{document}
